% ID: SEL-D0024-OB-E04
% ESTADO: cotejada. Segunda pasada visual completa y borrador compilado revisados.
% Fuente: Selectividad/matematicas_aplicadas/reserva_b_EXAMEN 2-ANDALUCÍA.pdf
% SHA-256: 98b7723232acd3143de28397f7641923ed48292db7c35883ebd2b56a72b30b79
% Página: 2; opción B; ejercicio original 4.

El tiempo diario, en horas, que dedican los alumnos de una Facultad a las redes sociales sigue una ley Normal de desviación típica 2 horas. Se toma una muestra aleatoria de 10 alumnos con los siguientes tiempos en horas
\[
6.5\quad7\quad6.25\quad7\quad5.5\quad7.25\quad6.75\quad6.25\quad6\quad6.5
\]
\begin{enumerate}[label=\alph*)]
\item \textbf{(1.5 puntos)} Determine el intervalo de confianza, al 90\%, para el tiempo medio diario dedicado por los alumnos de esa Facultad a las redes sociales.
\item \textbf{(1 punto)} Utilizando el mismo nivel de confianza anterior, calcule el tamaño muestral mínimo necesario para estimar el tiempo medio diario, para un error de estimación máximo de 0.1 horas.
\end{enumerate}
